% Estilo do PDF do atlas, no padrão visual dos documentos do Programa Youthmappers UFV.
% Não define cor nem fonte: usa só os nomes de identidade.tex, que construir.py gera a partir do
% _brand.yml (atlasprimaria, atlassecundaria, atlasdestaque, atlasclaro, atlastexto,
% \atlaslogos, \atlasmarca). Entra no preâmbulo pelo include-in-header do _quarto.yml.
\usepackage{fontspec}
\input{identidade.tex}
\usepackage{tikz}
\usepackage{enumitem}
\usepackage{caption}
\usepackage{scrlayer-scrpage}
\usepackage{colortbl}

\color{atlastexto}
\KOMAoptions{parskip=half}
\setlength{\emergencystretch}{3em}
\frenchspacing
\hyphenation{OpenStreetMap}

% Títulos: sem serifa, na cor principal, com filete na cor de destaque
\setkomafont{disposition}{\sffamily\bfseries\color{atlasprimaria}}
\addtokomafont{chapter}{\huge}
\addtokomafont{section}{\Large}
\addtokomafont{subsection}{\large}
\RedeclareSectionCommand[beforeskip=0pt, afterskip=1.2\baselineskip]{chapter}
\RedeclareSectionCommand[beforeskip=-1.2\baselineskip, afterskip=0.5\baselineskip]{section}

% "Parte" (Pranchas): faixa na cor principal acima do capítulo seguinte, sem página própria
\newif\ifatlasparte
\newcommand{\atlasparte}{}
\renewcommand{\part}[2][]{\clearpage\global\atlaspartetrue\gdef\atlasparte{#2}%
  \addcontentsline{toc}{part}{#2}}
\newcommand{\atlasfaixaparte}{%
  \ifatlasparte
    \colorbox{atlasprimaria}{\parbox{\dimexpr\linewidth-2\fboxsep}{%
      \vspace{3pt}\normalsize\sffamily\bfseries\color{white}\MakeUppercase{\atlasparte}\vspace{3pt}}}%
    \par\vspace{0.8\baselineskip}\global\atlaspartefalse
  \fi}
\renewcommand*{\chapterlinesformat}[3]{\atlasfaixaparte#2#3\par\vspace{-4pt}%
  {\color{atlasdestaque}\rule{\linewidth}{1.2pt}}}

% Cabeçalho (título do atlas e nome da marca) e número da página
\clearpairofpagestyles
\KOMAoptions{headsepline=0.4pt}
\setkomafont{headsepline}{\color{atlasdestaque}}
\addtokomafont{pagehead}{\sffamily\footnotesize\color{atlastexto!75}\upshape}
\addtokomafont{pagefoot}{\sffamily\footnotesize}
\ihead{\atlastitulo}
\ohead{\atlasmarca}
\cfoot*{\pagemark}
\pagestyle{scrheadings}
\renewcommand*{\chapterpagestyle}{scrheadings}

% Listas: marcador triangular na cor de destaque, números na cor secundária
\setlist{itemsep=2pt, topsep=3pt, parsep=0pt}
\setlist[itemize]{label=\textcolor{atlasdestaque}{\small$\blacktriangleright$}}
\setlist[enumerate]{label=\textcolor{atlassecundaria}{\sffamily\bfseries\arabic*.}}

% Tabelas: sem serifa, menores, filetes na cor principal
\AtBeginEnvironment{longtable}{\sffamily\small\arrayrulecolor{atlasprimaria}}
\renewcommand{\arraystretch}{1.25}

% Legendas, links e quadros (callouts do Quarto)
\captionsetup{font={small,sf}, labelfont={bf,color=atlasprimaria}}
\AtBeginDocument{\hypersetup{colorlinks=true, linkcolor=atlasprimaria, urlcolor=atlasprimaria,
  citecolor=atlasprimaria}}
\AtBeginDocument{%   o Quarto define as cores dos quadros depois deste preâmbulo
  \colorlet{quarto-callout-note-color}{atlasprimaria}%
  \colorlet{quarto-callout-tip-color}{atlasdestaque}%
  \colorlet{quarto-callout-warning-color}{atlassecundaria}%
  \colorlet{quarto-callout-important-color}{atlassecundaria}%
  \colorlet{quarto-callout-caution-color}{atlassecundaria}}

% Sumário
\renewcommand*{\contentsname}{Sumário}
\setkomafont{chapterentry}{\sffamily\bfseries\color{atlasprimaria}}
\setkomafont{partentry}{\sffamily\bfseries\color{atlassecundaria}}
